% NORMA_COPPE_2026.tex -- LaTeX source for the typeset PDF version of the
% proposed minimal COPPE/UFRJ norm. Uses the article class on purpose:
% this is a normative document about the coppe class, not a thesis, so
% applying the coppe class itself to it would be confusing.
%
% Authoring source is NORMA_COPPE_2026.md at the repository root; this
% file is a hand-typeset transcription kept readable as LaTeX without
% pandoc.

\documentclass[a4paper,11pt]{article}
\usepackage[T1]{fontenc}
\usepackage[utf8]{inputenc}
\usepackage{lmodern}
\usepackage[brazilian]{babel}
\usepackage[margin=2.5cm]{geometry}
\usepackage{titlesec}
\usepackage{enumitem}
\usepackage{microtype}
\usepackage{url}
\usepackage{parskip}
\usepackage[hidelinks]{hyperref}

\titleformat{\section}{\large\bfseries}{\thesection}{1ex}{}
\titleformat{\subsection}{\normalsize\bfseries}{\thesubsection}{1ex}{}

\title{\bfseries Norma para a Elabora\c{c}\~ao Gr\'afica de Teses e
       Disserta\c{c}\~oes da COPPE/UFRJ\\[2mm]
       \large Edi\c{c}\~ao 2026 -- Proposta para aprecia\c{c}\~ao da CPGP}
\author{}
\date{Maio de 2026}

\begin{document}
\maketitle

\noindent\textit{Esta Norma, quando aprovada, substitui a} Norma para a
Elabora\c{c}\~ao Gr\'afica de Teses/Disserta\c{c}\~oes \textit{da
COPPE/UFRJ aprovada pela CPGP em 15 de julho de 2008 e suas revis\~oes
subsequentes (01/10/2009, 10/09/2010 e 26/11/2019).}

\hrulefill

\section{Documento base e implementa\c{c}\~ao de refer\^encia}

\subsection{Documento base}

A apresenta\c{c}\~ao gr\'afica de teses e disserta\c{c}\~oes da
COPPE/UFRJ segue o \textbf{Manual para Elabora\c{c}\~ao e
Normaliza\c{c}\~ao de Trabalhos Acad\^emicos} da UFRJ/SiBI, 9.\textsuperscript{a}
edi\c{c}\~ao (Rio de Janeiro, 2025) -- doravante referido como
\textit{Manual UFRJ 2025} -- em conjunto com as adapta\c{c}\~oes
estabelecidas pela presente Norma.

A presente Norma estabelece \textbf{apenas os elementos pr\'oprios da
COPPE/UFRJ}. Todos os aspectos n\~ao tratados explicitamente aqui --
formato do papel, margens, espa\c{c}amento, fonte, pagina\c{c}\~ao
geral, equa\c{c}\~oes, tabelas, indicativos de se\c{c}\~ao, ortografia,
siglas, ilustra\c{c}\~oes no caso geral, estrutura textual no caso
geral, etc.\ -- devem seguir o \textit{Manual UFRJ 2025}.

\subsection{Implementa\c{c}\~ao de refer\^encia}

A \textbf{implementa\c{c}\~ao de refer\^encia} desta Norma \'e a classe
LaTeX \texttt{coppe}, vers\~ao 4.0 e posteriores (\textbf{CoppeTeX
4.x}), distribu\'ida em
\url{https://github.com/COPPE-UFRJ/CoppeTeX}. A classe produz, sem
configura\c{c}\~ao adicional por parte do autor, um documento aderente
simultaneamente ao \textit{Manual UFRJ 2025} e \`a presente Norma. Seu
uso \'e o caminho recomendado pela COPPE/UFRJ para a confec\c{c}\~ao do
documento final.

Quando houver d\'uvida quanto \`a apresenta\c{c}\~ao correta de algum
elemento gr\'afico previsto pela presente Norma, a sa\'ida gerada pela
vers\~ao mais recente do CoppeTeX 4.x prevalece como refer\^encia.

\section{Identifica\c{c}\~ao institucional}

\subsection{Capa e folha de rosto}

A capa e a folha de rosto do trabalho exibem, na ordem:

\begin{itemize}[leftmargin=2em]
  \item \textbf{Universidade Federal do Rio de Janeiro};
  \item \textbf{Instituto Alberto Luiz Coimbra de P\'os-Gradua\c{c}\~ao
        e Pesquisa de Engenharia (COPPE)};
  \item O nome do \textbf{Programa de P\'os-gradua\c{c}\~ao} em que o
        trabalho foi defendido;
  \item O nome completo do autor;
  \item O t\'itulo do trabalho, no idioma principal (ver Se\c{c}\~ao 4);
  \item A cidade (Rio de Janeiro) e o ano da defesa.
\end{itemize}

\subsection{Car\'ater da identidade institucional}

Os elementos de identidade institucional -- o nome da Universidade, o
nome do Instituto (COPPE), o nome do Programa de P\'os-gradua\c{c}\~ao,
a cidade sede (Rio de Janeiro), o estado (RJ) e o pa\'is (Brasil) --
\textbf{s\~ao sempre exibidos em portugu\^es}, independentemente do
idioma principal do trabalho. Somente o t\'itulo do trabalho, na capa e
na folha de rosto, acompanha o idioma principal.

\subsection{Ficha catalogr\'afica}

A ficha catalogr\'afica \'e gerada conforme as regras AACR2 e o
\textit{Manual UFRJ 2025}, e figura no \textbf{verso da folha de rosto}.
Acrescenta-se \`a ficha:

\begin{itemize}[leftmargin=2em]
  \item A indica\c{c}\~ao da institui\c{c}\~ao como \textbf{UFRJ/COPPE};
  \item O nome do \textbf{Programa de P\'os-gradua\c{c}\~ao} em que o
        trabalho foi defendido;
  \item A indica\c{c}\~ao das p\'aginas iniciais e finais da lista de
        Refer\^encias.
\end{itemize}

\subsection{Folha de aprova\c{c}\~ao}

A folha de aprova\c{c}\~ao cont\'em o t\'itulo do trabalho, o nome do
candidato, os nomes do(s) orientador(es) e dos membros da Banca
Examinadora -- na ordem em que aprovaram o trabalho --, com espa\c{c}o
para assinatura, e a data da defesa. A folha de aprova\c{c}\~ao
\textbf{n\~ao \'e numerada}.

\section{Resumos: estrutura de tr\^es idiomas}

Em complemento ao previsto no \textit{Manual UFRJ 2025}, a presente
Norma reconhece tr\^es posi\c{c}\~oes distintas para o resumo:

\begin{enumerate}[leftmargin=2em]
  \item \textbf{Resumo em idioma principal}, obrigat\'orio, no idioma
        em que o trabalho foi redigido (ver Se\c{c}\~ao 4);
  \item \textbf{Resumo em idioma estrangeiro} (\textit{foreign
        abstract}), obrigat\'orio, por conven\c{c}\~ao em ingl\^es;
        quando o idioma principal \'e o ingl\^es, este resumo \'e
        redigido em portugu\^es;
  \item \textbf{Resumo em portugu\^es}, opcional e adicional, requerido
        apenas quando nenhum dos dois resumos anteriores seja em
        portugu\^es -- caso t\'ipico de trabalhos com idioma principal
        espanhol, franc\^es ou italiano. Destina-se \`a leitura pela
        banca e pelos examinadores brasileiros.
\end{enumerate}

Cada resumo ocupa uma p\'agina pr\'opria, na ordem acima, dentro da
se\c{c}\~ao pr\'e-textual.

\section{Idioma principal do trabalho}

O trabalho pode ser redigido em qualquer um dos seguintes idiomas
principais:

\begin{itemize}[leftmargin=2em]
  \item \textbf{Portugu\^es};
  \item \textbf{Ingl\^es};
  \item \textbf{Espanhol};
  \item \textbf{Franc\^es};
  \item \textbf{Italiano};
  \item Qualquer outro idioma reconhecido pela implementa\c{c}\~ao de
        refer\^encia (CoppeTeX 4.x), mediante o pacote de
        localiza\c{c}\~ao correspondente.
\end{itemize}

A escolha do idioma principal afeta a l\'ingua do corpo textual, das
legendas, dos t\'itulos das se\c{c}\~oes pr\'e-textuais (Sum\'ario,
Listas, Refer\^encias, etc.) e do t\'itulo do trabalho exibido na capa.
A identidade institucional permanece em portugu\^es em qualquer idioma
principal (ver Subse\c{c}\~ao 2.2).

\section{Ilustra\c{c}\~oes pr\'oprias da COPPE/UFRJ}

Em complemento \`as listas previstas no \textit{Manual UFRJ 2025}
(Lista de Figuras, Lista de Tabelas, Lista de Abreviaturas e Siglas,
Lista de S\'imbolos), a presente Norma reconhece:

\begin{itemize}[leftmargin=2em]
  \item \textbf{Lista de Quadros} -- para o tipo de ilustra\c{c}\~ao
        ``quadro'', bordado em todos os lados, distinto de tabela (a
        qual \'e delimitada apenas no topo e na base);
  \item \textbf{Lista de Programas} -- para listagens de
        c\'odigo-fonte;
  \item \textbf{Lista de Algoritmos} -- para a apresenta\c{c}\~ao de
        pseudo-c\'odigo.
\end{itemize}

Cada uma destas listas figura entre a Lista de Tabelas e a Lista de
Abreviaturas, na se\c{c}\~ao pr\'e-textual, sempre que houver ao menos
uma ocorr\^encia no corpo do trabalho.

Para todas as ilustra\c{c}\~oes (figuras, tabelas, quadros, programas,
algoritmos) a legenda figura \textbf{acima} da ilustra\c{c}\~ao, e a
fonte \'e obrigat\'oria e figura \textbf{abaixo}, em corpo menor,
espa\c{c}amento simples, centralizada, antecedida da palavra
``Fonte:''\ (ou seu equivalente no idioma principal).

\section{Ap\^endices e Anexos}

Os cap\'itulos de Ap\^endice e de Anexo s\~ao apresentados, tanto no
in\'icio do cap\'itulo quanto no Sum\'ario, no formato:

\begin{quote}
  \textbf{AP\^ENDICE A -- \textit{T\'itulo do ap\^endice}}\\[1mm]
  \textbf{ANEXO A -- \textit{T\'itulo do anexo}}
\end{quote}

A palavra ``AP\^ENDICE''\ ou ``ANEXO''\ figura em caixa alta, seguida
da letra identificadora (A, B, C, \ldots), de travess\~ao e do
t\'itulo do cap\'itulo, respeitando o idioma principal do trabalho.

\section{Numera\c{c}\~ao das p\'aginas pr\'e-textuais}

A capa n\~ao \'e contada nem numerada. A folha de rosto inicia a
contagem das p\'aginas pr\'e-textuais.

Por \textbf{regra padr\~ao}, \textbf{os n\'umeros das p\'aginas
pr\'e-textuais n\~ao s\~ao impressos}, embora sejam contados. \`A
crit\'erio do autor, e de forma opcional, os n\'umeros das p\'aginas
pr\'e-textuais podem ser impressos em algarismos romanos min\'usculos
(iii, iv, v, \ldots) a partir da terceira p\'agina contada.

A numera\c{c}\~ao ar\'abica, no canto superior externo da p\'agina,
inicia em 1 na primeira p\'agina da parte textual.

\section{Cita\c{c}\~oes e Refer\^encias Bibliogr\'aficas}

A presente Norma adota integralmente as edi\c{c}\~oes mais recentes
das normas ABNT pertinentes:

\begin{itemize}[leftmargin=2em]
  \item \textbf{NBR 6023} (Refer\^encias) em sua revis\~ao a partir de
        2018;
  \item \textbf{NBR 10520} (Cita\c{c}\~oes) em sua revis\~ao de 2023.
\end{itemize}

Os elementos da refer\^encia, os formatos das cita\c{c}\~oes e as
regras gerais de apresenta\c{c}\~ao s\~ao os definidos pela ABNT
vigente.

Esta ado\c{c}\~ao constitui a \'unica diverg\^encia da presente Norma
em rela\c{c}\~ao ao previsto no \textit{Manual UFRJ 2025} -- que
reflete edi\c{c}\~oes anteriores das mesmas normas ABNT -- e ocorre na
medida estrita das atualiza\c{c}\~oes posteriores a 2020 introduzidas
pela pr\'opria ABNT.

A implementa\c{c}\~ao de refer\^encia (CoppeTeX 4.x) fornece
automaticamente a formata\c{c}\~ao correta para todos os tipos de
documento previstos na NBR 6023 vigente, em qualquer dos sistemas de
chamada (autor-data ou num\'erico) \`a escolha do autor.

\section{Implementa\c{c}\~ao de refer\^encia}

A classe LaTeX \texttt{coppe} em suas vers\~oes 4.0 e posteriores
(\textbf{CoppeTeX 4.x}), distribu\'ida em
\url{https://github.com/COPPE-UFRJ/CoppeTeX}, \'e a implementa\c{c}\~ao
de refer\^encia desta Norma. Sua ado\c{c}\~ao \'e o caminho recomendado
pela COPPE/UFRJ para garantir ader\^encia autom\'atica a todos os
requisitos acima e ao \textit{Manual UFRJ 2025}.

Em caso de diverg\^encia entre a presente Norma escrita e a sa\'ida
gerada pela vers\~ao est\'avel corrente do CoppeTeX 4.x, prevalece a
presente Norma e o \textit{Manual UFRJ 2025}, e a diverg\^encia deve
ser comunicada \`a equipe mantenedora do CoppeTeX para
corre\c{c}\~ao na pr\'oxima revis\~ao da implementa\c{c}\~ao.

\vspace{4mm}\hrulefill

\noindent\textit{Vers\~ao da proposta: maio de 2026. \emph{Branch}
\texttt{nlinguas} em
\url{https://github.com/COPPE-UFRJ/CoppeTeX}.}

\end{document}
